% !TeX root = ../papers/nc-p3.tex
\section{指数二の安定化点の孤立性と八次曲面上の剛な四束}
\label{sec:octic-surface-fixed-center}
\begin{theorem}\label{thm:octic-surface-isolated-inertia}
十五次元閉包の有限四辺にある正則整域窓を取り、圧縮次数の中心を $Z$、$X=\operatorname{Proj} Z$、
反射的偏極を $H$ とする。位数二の次数安定化群を持つ射影点は高々有限個である。
一般 $S\in|2H|$ は安定化点を避け、滑らか、または有限個の $A_1$ 点を持つ。
$h=H|_S$ は豊富な直線束で $K_S=h$, $h^2=8$, $p_g=5$, $q=0$ である。
\end{theorem}
\begin{proof}
系\ref{cor:index-four-no-divisorial-inertia}で因子安定化は除かれている。局所表示を明示する。
零にならない次数2の切断 $b$ で、錐の $b=1$ の切片を $Y$ とし、
$Y$ 上の符号作用で割ると $X$ のこのチャートを得る。
折り畳みも $i$ から $j$ の経路長 $4m+j-i$ に次数 $m$ を付けて同じ切片を取る。
この次数作用は元の経路長の倍率と頂点基底変更の組であり、単純の置換は
定理\ref{thm:central-character-periods}の圧縮次数作用と同じである。
$b$ を逆にしてその平方根をエタールに添加した錐は $Y\times\mathbb G_m$ となり、
代数も対応するLaurent拡大である。従って切片の代数は有限大域次元、極大Cohen--Macaulay性と
中心trace双対を持つ。大域次元についてはLaurent拡大した加群の分解を制限し、
元の加群を直和因子に取ってもよい。 $Y$ はGorensteinで、符号作用は標準形式の生成元に $-1$ を掛ける。
実際 $\omega_Z=Z(-1)$ の斉次生成元は奇次数であり、$b-1$ は不変な正則方程式なので
随伴で同じ指標を取る。

もし安定化曲線があれば、基点のない豊富な $|2H|$ の幾何学的一般切断でその一般部を切る。
切片 $Y$ 上にもこの切断を引き戻す。普遍接続が射影空間束であるため、接続での多項式拡大・
局所化から一般曲面上の代数も有限大域次元となる。
定理\ref{thm:central-cut-transverse-a-one}の局所論証をここで適用できる。
すなわち曲面の中心は正規Gorensteinで、極大Cohen--Macaulayな代数はtrace自己双対、
有限大域次元によりhomologically homogeneousで中心は有理特異点を持つ。
従って特異ならDu Valで、穴あき完備曲面のBrauer群が零であることと反射的延長から
代数は $\operatorname{End}_R(M)$ になる。全直既約反射的加群が $M$ に現れる。
薄い単純の有限個の指標は切断が避けるので、各単純の頂点次元は2以上であり、
一頂点の総階数4により特異なら $A_1$、二種類の単純の頂点次元はいずれも2である。
一般曲面上の不変な局所切断方程式で随伴を取るので、標準形式の指標は再び $-1$ である。
曲面上の固定点は孤立する。元に固定因子がなく、一般切断はどの固定曲線も含まないからである。

固定点がこの曲面の滑らかな点なら接空間の二固有値はともに $-1$ である。
有限群は完備正則環のパラメータを平均化して線型化でき、固有値1があれば固定曲線ができる。
しかし二つが $-1$ なら標準形式上の作用は $+1$ で矛盾する。
残る $A_1$ 点で $R$ を幾何的完備局所環、$R'$ を符号作用の不変環とする。
穴あき部分の作用は自由で、そこで代数は商の穴あき曲面上のAzumaya環に降下する。
この粗商の穴あき曲面のBrauer群も零である。この事実に必要な局所仮定を確認する。
$R=k[[u,v]]^{\{\pm1\}}$ の穴あき部分の二重被覆は普遍エタール被覆である
（正則厳密局所環の純性）。従って符号作用はこの被覆に持ち上がり、正規化で原点まで延びる。
全持ち上げは位数4の有限群で、$R'$ はその正則曲面の商である。
有限群の線型化と曲面商特異点の解消により例外は有理曲線の木であり、
先の横断型の証明と同じ剰余・純性・形式存在の議論でBrauer群は零になる。
ここで剰余体は幾何学的一般切断の代数閉体である。

よって粗商上の分裂束を穴あき $R$ へ引き戻し、反射的に延長して
\emph{符号同変な} $M$ を取れる。その自己準同型環と元の代数の同定も同変である
（穴あき部分で得た同定の一意な反射的延長）。$A_1$ の直既約反射的加群の二型は
自由な $R$ と非自由な階数一の $I$ なので、任意の底の自己同型は両型を別々に保つ。
従って同変な $M$ から得る作用は自己準同型環の二つの単純型をそれぞれ固定する。
実際Krull--Schmidt分解の各同型型が半単純商の一つの中心冪等元に対応する。
しかし次数作用で固定される非薄単純には頂点次元が少なくとも $2\cdot2=4$ 必要である
（定理\ref{thm:central-character-periods}の固有層分解）。ここでは二型とも頂点次元2で矛盾する。
これで曲線上の安定化は除かれた。

一般切断は残る有限個の安定化点を避ける。そこで $H$ は直線束で、
曲面の局所型は先の切片での議論による。$K_X=-H$ と随伴から $K_S=h$、$H^3=4$ より $h^2=8$。
中心は深さ4で、切断完全列と $\dim Z_1=5$ から $h^0(h)=5$, $H^1(\mathcal O_S)=0$ を得る。
\end{proof}

\begin{theorem}\label{thm:octic-surface-fixed-center}
同じ仮定で、実現された一つの次数付き中心を固定すると、その上の窓の\emph{添字保持同型類}は有限である。
四次中心系に基点があってもよい。幾何学的一般八次曲面の標準滑らかなスタック上で、
一つの $\mu_4$-gerbe、四つのねじれ階数4束 $E_i$ と $h$ から元の窓を復元できる。
$H^0(\mathcal End E_i)=k'$, $H^1(\mathcal End E_i)=0$, $h^2(\mathcal End E_i)=35$ である。
ここで $k'$ はこの幾何学的一般切断の係数体で、最後の $h^2$ はコホモロジー次元である。
\end{theorem}
\begin{proof}
曲面の滑らかな部分上で零番隅の行列束の $\mathrm{SL}_4$ 持上げgerbeを取る。
他の頂点をつなぐ隅は森田同値を与え、共通のgerbe上に四束が定まる。
曲面の標準スタックは $A_1$ の位置で $[\mathbb A^2/\{\pm1\}]$ の局所型を持つ。
純性で有限係数gerbeの類を延ばし、反射的に延長した四束は滑らかな曲面上で局所自由である。
そのHomを粗曲面へ押したものは元の隅の制限である。穴あき部分で一致し、いずれも反射的である。
有限安定群の不変部分は完全なのでコホモロジーはこれで計算できる。
$E_{i+4}=E_i(h)$ と置く。中心八次切断の完全列と隅の深さ4により、
経路長8未満の $H^0(\mathcal Hom(E_i,E_j))$ は経路空間そのもので、合成も保たれる。
一・二文字とその積から全ての二次関係が戻る。特に零番からの四辺を区別した同型類を復元する。
自己Homの第一コホモロジーは同じ完全列の両側の中間コホモロジー消滅から零、
零次は定数であり、Serre双対と $K_S=h$ により二次の次元は経路長4の次元35である。

有限性のため、まず固定中心 $Z$ を印した候補が有限型の係数空間で覆えることを示す。
$Z$ の有限斉次表示とその四つの八次パラメータを一度固定する。四窓テンソルの
有限Frobenius・有限辺条件の上で、各中心生成元の像をその長さの経路空間に選ぶ。
有限次元の表示行列の階数別に基底を取れる。中心性と定義関係を課し、指定した四パラメータで
定理\ref{thm:index-four-octic-normalization}の有限検査を課す。十五次元閉包内に限り、これは有限型の
構成可能な条件で、正則整域を与える。その四パラメータ上で有限かつ次元4なのでパラメータの像は代数的独立、
$Z$ はその上の有限整域だから写像 $Z\to Z(B)$ は単射である。この閉包の中心Hilbert関数は一定だから全射でもある。
逆に固定中心を持つ全窓がこの表示に現れる。

この有限型性から、以下の剛性を幾何学的一般点で使えば十分である。
係数空間の任意の既約部分族 $T$ を取り、$T$ のパラメータとは独立に固定中心の
幾何学的一般切断を取り、その体 $k'$ に底を変える。曲面とその標準スタックはもはや固定される。
$T$ の幾何学的一般点で得たgerbeの類はこの固定スタック上の類から来る。
代数閉体の拡大で有限係数のエタールコホモロジーは変わらないからである（有限安定群のチャートでも同じ）。
その一つのgerbe上の四束と必要な同定は有限表示なので、$T_{k'}$ の有限な一般体拡大の後、
稠密開集合に族として広げられる。復元した一・二文字空間・乗法と元の族の同定も同時に広げられる。

固定した固有スタック上の束の族で各自己第一Extが零なら、一般点の近くで各束の同型型は一定となる。
詳しくは開集合の一般な幾何閉点で束 $E$ を固定し、族 $E_t$ に対する $\mathcal Hom(E,E_t)$ の
第一コホモロジーがその点で零であることを使う。コホモロジーと底変換により恒等写像を近くの族の射に延ばせる。
底変換は粗空間への完全な直像と固有射で適用できる。行列式が可逆なことは固有性で底上開条件だから
近くで同型になる。四束に同時に使うと、復元により元の窓の幾何的同型型も稠密開集合上で一つである。
四窓テンソルの基底変更軌道による同型は有限型の条件なので、代数閉体拡大で新たな同型は生じない。
例えば軌道関係の逆像が $T\times T$ の一般点を含むことを上の独立な底拡大で検査し、
構成可能性から $T$ の一般な閉点の同型類が稠密開集合を含むと分かる。
補集合にNoether帰納を適用して有限性を得る。ここでは中心の全ての印付け自体の個数は数えていない。
\end{proof}
